\documentclass[11pt]{article}

\usepackage[utf8]{inputenc}
\usepackage[T1]{fontenc}
\usepackage{geometry}
\usepackage{amsmath}
\usepackage{amssymb}
\usepackage{booktabs}
\usepackage{microtype}
\usepackage{xcolor}
\usepackage{hyperref}

\geometry{margin=1in}
\hypersetup{
  colorlinks=true,
  linkcolor=blue!50!black,
  citecolor=blue!50!black,
  urlcolor=blue!50!black
}

\title{LitJev: A Calibrated Single-Pass Decision Model for Academic Manuscript Screening}
\author{\textbf{Replace with author metadata before submission}}
\date{arXiv preprint, October 2026}

\begin{document}

\maketitle

\begin{abstract}
Evaluating academic manuscripts with large generative models is expensive,
slow, and produces free-form assessments that are difficult for software
agents to consume. This paper presents LitJev, a 4B-parameter single-pass
decision model that maps a manuscript title and abstract to twelve typed
decisions, each with a calibrated probability, without autoregressive
generation or growing KV-cache. The model is trained with three groups of
supervision: craft questions use teacher distillation as a cold-start signal,
value questions use long-term academic outcomes and knowledge-manifold
trajectory labels rather than current large-model consensus, and
integrity/coverage questions use retraction, manipulation-cluster,
programmatic-perturbation, and out-of-domain signals. On a frozen,
year-held-out in-domain benchmark, eight of twelve question heads achieve AUROC
above 0.95, eleven of twelve achieve post-hoc expected calibration error at or
below 0.026, and the explicit domain-coverage head reaches AUROC 0.996. An
exploratory comparison shows that the specialized model produces lower Brier
score than zero-shot general-purpose models on the same typed propositions,
while several rare-class value tasks remain limited by supervision. We
intentionally do not optimize the long-term-impact rank-correlation beyond a
regression gate, because citation alignment is not the same as evidence that
the model can recognize silent frontiers.
\end{abstract}

\section{Introduction}

Scientific publishing increasingly depends on preliminary screening under
resource constraints. A complete expert review is expensive and slow, while a
generative LLM review is faster but has several limitations for
machine-to-machine workflows: it is costly at scale, it emits free-form prose
that downstream systems must parse, it is often uncertain in ways that are not
expressed as usable probabilities, and, if used as a value judge, it tends to
reproduce current mainstream consensus.

LitJev is a deliberate System-1 counterpart to generative System-2 review. It
is not a reviewer, a chatbot, or a journal decision system. It is a
task-specific decision model that produces a fixed vector of per-question
probabilities in one forward pass. The contributions of this work are:

\begin{enumerate}
  \item a twelve-question typed decision schema spanning craft, value,
        integrity, and domain-coverage dimensions;
  \item a three-signal supervision design in which only time-invariant craft
        dimensions are teacher-distilled, while value dimensions use academic
        outcomes and trajectory geometry instead of current LLM opinion;
  \item a 4B-parameter cross-encoder implementation with soft-label training
        and a pair-ranking objective for long-term-impact thresholds;
  \item a frozen in-domain benchmark with explicit calibration and OOD
        abstention;
  \item an exploratory zero-shot comparison with two general-purpose models,
        framed as evidence about task specialization rather than a superiority
        claim.
\end{enumerate}

\section{Related Work}

\paragraph{System-1 decision models.}
Jev introduced the idea of training models to produce typed decisions and
calibrated probabilities instead of generating text. OpenJev and mini-jev
showed that compact cross-encoders can approximate this downstream shape for
benchmark-style inference. RLCD has not been fully publicly disclosed; we
therefore follow RLCR only as a conceptual reference and do not claim to
reproduce RLCD.

\paragraph{Outcome-based review.}
ReviewGuard aligns review signals toward future citation impact. APRES searches
for rubrics that predict future citations. These systems demonstrate that
long-term outcome signals can be learned, but they also measure
consensus-compatible ranking signals. Our design keeps these signals for
instrumented monitoring but does not make them the primary target of all value
supervision.

\paragraph{Knowledge manifold learning and silent frontiers.}
KML models citation trajectories as a graph manifold and defines source nodes,
silent-frontier trajectories, and manipulation clusters. This family of signals
is used for value and integrity supervision because it explicitly addresses
papers whose early citation patterns do not yet reflect long-term influence.

\paragraph{Calibration and abstention.}
Temperature scaling and isotonic regression are standard post-hoc calibration
tools. Selective prediction and out-of-domain detection are relevant for models
that should abstain outside their training community. We implement an
\texttt{in\_coverage} question head rather than a generative refusal.

\section{Task Formulation}

LitJev receives a structured paper state. The current snapshot uses only title
and abstract. The output is a vector
\(
  \mathbf{y}=[p_1,\dots,p_{12}],\quad p_i\in[0,1].
\)

The twelve question heads are grouped as shown in
Table~\ref{tab:schema}.

\begin{table}[t]
\centering
\begin{tabular}{@{}lll@{}}
\toprule
Group & Question & Output \\
\midrule
Craft & \texttt{clarity} & choice posterior \\
Craft & \texttt{method\_type} & choice posterior \\
Craft & \texttt{rigor} & 1--5 threshold set \\
Craft & \texttt{reproducibility} & 1--5 threshold set \\
Craft & \texttt{experiments\_support} & predicate probability \\
Value & \texttt{long\_term\_impact} & 1--5 threshold set \\
Value & \texttt{expert\_accept} & predicate probability \\
Value & \texttt{silent\_source\_potential} & predicate probability \\
Value & \texttt{paradigm\_trajectory} & five-class posterior \\
Integrity & \texttt{citation\_chain\_clean} & predicate probability \\
Integrity & \texttt{misconduct\_signal} & predicate probability \\
Coverage & \texttt{in\_coverage} & predicate probability \\
\bottomrule
\end{tabular}
\caption{Question heads of the LitJev decision schema.}
\label{tab:schema}
\end{table}

The model does not output a total score and does not rank papers against one
another. Each probability is exposed independently so that a downstream agent
can enforce its own decision rules.

\section{Supervision Design}

\subsection{Craft dimensions}

Clarity, method type, rigor, reproducibility, and experimental support are
treated as time-invariant properties. They are cheap for a strong teacher to
estimate and are not the main source of value bias. We produce soft labels by
prompting a commercial teacher with five reviewer personas per paper and taking
the vote distribution.

\subsection{Value dimensions}

Value labels are not teacher-distilled. \texttt{long\_term\_impact} is derived
from citation counts normalized within venue--year cohorts.
\texttt{expert\_accept} is a human decision anchor.
\texttt{silent\_source\_potential} and \texttt{paradigm\_trajectory} use KML
source and trajectory criteria.

The reason is conceptual: frontier ideas can be systematically misestimated by
consensus models during the silent period. Distilling a general-purpose LLM for
these questions would therefore train on the wrong direction for a key
differentiator.

\subsection{Integrity and coverage}

Integrity labels combine retraction status, KML manipulation clusters, and
programmatically perturbed negatives. The coverage question is trained with
in-domain ML-community papers as positives and twelve OpenAlex scientific
domains as negatives. It allows the model to express ``this input is outside my
coverage'' as a first-class decision.

\section{Data}

\subsection{Sources}

The main community corpus contains ICLR and NeurIPS submissions with decision
metadata, augmented with OpenAlex citation information, Semantic Scholar
matching, arXiv identifiers, and KML trajectory features. OOD negatives are
sampled from mathematics, physics, chemistry, biology, medicine, economics,
psychology, geology, materials science, sociology, history, and philosophy.

\subsection{Pair construction}

Each decision sample becomes one or more NLI-style premise--hypothesis pairs.
Choice questions emit one pair per value. Score questions emit threshold
propositions (``at least $k$''). Binary predicates emit one proposition. A paper
is the premise; a proposition is the hypothesis.

\subsection{Final snapshot}

\begin{table}[t]
\centering
\begin{tabular}{@{}lc@{}}
\toprule
Split & Pair count \\
\midrule
train & 268{,}039 \\
validation & 137{,}987 \\
test & 109{,}607 \\
\midrule
total & 515{,}633 \\
\bottomrule
\end{tabular}
\caption{Final supervised snapshot.}
\label{tab:data}
\end{table}

The test split is held out by year and is not used for model selection or
calibration fitting. Calibration mappings are fit on validation only.

\section{Model and Training}

The model uses Qwen3.5-4B as a decoder trunk, takes the last token of the
concatenated input, and maps the pooled representation to a two-way head. The
positive-class logit is softmax-normalized into the decision probability.

Training uses a soft cross-entropy loss against target probabilities and a
RankNet-style pair loss for long-term-impact papers within the same
venue--year cohort,
\(
  \mathcal{L}=\mathcal{L}_{\mathrm{CE}}+\lambda\mathcal{L}_{\mathrm{rank}}.
\)
The full parameter set is fine-tuned in bf16 with gradient checkpointing.
Sortish length bucketing reduces padding. Checkpoints store optimizer and
scheduler state for exact resumption.

After training, each question is calibrated independently on validation using
temperature scaling and isotonic regression. Final head outputs are
monotonically remapped before use.

\section{Evaluation}

\subsection{Metrics}

\begin{itemize}
  \item \textbf{AUROC}: binary discrimination for predicate and threshold
        questions.
  \item \textbf{Brier}: squared error against the stored target probability.
  \item \textbf{ECE}: ten-bin expected calibration error after post-hoc
        calibration.
  \item \(\rho_{\mathrm{LTI}}\): cohort-level Spearman correlation of per-paper
        expected score.
  \item \textbf{high\_impact\_rejected\_recall}: recovered high-impact rejected
        manuscripts.
\end{itemize}

\subsection{Comparators}

For an exploratory comparison, we evaluate GLM-5.3 and GLM-5.3-Flash on the
same 192-pair stratified test sample with zero-shot prompting and JSON-encoded
probabilities. The models are general-purpose and are not calibrated; the
comparison is not intended to establish general capability.

\section{Results}

\subsection{Frozen held-out benchmark}

\begin{table}[t]
\centering
\begin{tabular}{@{}lcc@{}}
\toprule
Question & AUROC & \(n\) \\
\midrule
\texttt{citation\_chain\_clean} & 1.0000 & 1{,}590 \\
\texttt{in\_coverage} & 0.9963 & 2{,}587 \\
\texttt{method\_type} & 0.9893 & 24{,}670 \\
\texttt{rigor} & 0.9887 & 19{,}736 \\
\texttt{clarity} & 0.9803 & 14{,}802 \\
\texttt{reproducibility} & 0.9756 & 19{,}736 \\
\texttt{misconduct\_signal} & 0.9755 & 2{,}277 \\
\texttt{paradigm\_trajectory} & 0.8780 & 8{,}235 \\
\texttt{experiments\_support} & 0.7732 & 4{,}934 \\
\texttt{silent\_source\_potential} & 0.7607 & 974 \\
\texttt{expert\_accept} & 0.6966 & 3{,}794 \\
\bottomrule
\end{tabular}
\caption{Frozen held-out test AUROC for eleven question heads.}
\label{tab:auc}
\end{table}

\texttt{long\_term\_impact} reports per-paper rank calibration rather than a
single AUROC. On the frozen test cohort,
\(\rho_{\mathrm{LTI}}=0.4707\) with \(n=1{,}568\). We present this number as a
regression baseline, not as the headline value capability.

\subsection{Calibration}

After isotonic calibration, ECE is at or below \(0.0261\) for eleven of twelve
question heads. The exception is \texttt{silent\_source\_potential}, with ECE
\(0.0595\).

\subsection{Training lineage}

\begin{table}[t]
\centering
\begin{tabular}{@{}lccc@{}}
\toprule
Snapshot & Paradigm test & Silent source test & \(\rho_{\mathrm{LTI}}\) test \\
\midrule
v2 & 0.913 & 0.757 & 0.474 \\
v3 & 0.857 & 0.762 & 0.470 \\
v3.1 & 0.878 & 0.761 & 0.470 \\
\bottomrule
\end{tabular}
\caption{Training lineage on the frozen test split.}
\label{tab:lineage}
\end{table}

The v3.1 checkpoint recovers most of the paradigm-trajectory regression while
preserving the silent-source value signal.

\subsection{Exploratory GLM comparison}

\begin{table}[t]
\centering
\begin{tabular}{@{}lccc@{}}
\toprule
Model & Brier & Hard accuracy & MAE \\
\midrule
GLM-5.3 & 0.1475 & 0.7558 & 0.2736 \\
GLM-5.3-Flash & 0.1578 & 0.7151 & 0.2956 \\
LitJev v3.1 & 0.0766 & 0.8663 & 0.1298 \\
\bottomrule
\end{tabular}
\caption{Paired zero-shot comparison, \(n=172\) valid pairs.}
\label{tab:glm}
\end{table}

The comparison is at most preliminary evidence for task specialization and
must not be read as a general-capability ranking.

\section{Discussion}

The main architectural value of LitJev is not a new model family but an
explicit separation between signals. Craft dimensions may be teacher-supervised
because they are time-invariant. Value dimensions require slow, structurally
defined signals because consensus models are biased in the exact regime that a
silent-frontier detector must handle.

We deliberately avoid optimizing \(\rho_{\mathrm{LTI}}\) beyond a regression
gate. Citation alignment is a useful instrumentation channel, but it is not
identical to recognizing under-appreciated work, and optimizing it too strongly
can push the model back toward mainstream consensus.

\section{Limitations}

\begin{itemize}
  \item The model uses title and abstract only; full-text understanding is
        future work.
  \item The community benchmark is ICLR/NeurIPS-centric and is not a general
        academic sample.
  \item \texttt{silent\_source\_potential} and
        \texttt{high\_impact\_rejected\_recall} have insufficient positive
        supervision for decisive conclusions.
  \item The GLM comparison is zero-shot, uncalibrated, and based on 172 paired
        examples.
  \item No serving-level latency or throughput benchmark is included.
  \item The study does not perform multi-seed replications, bootstrap tests, or
        a formal train/test contamination audit.
  \item Coverage is not a substitute for robustness outside the trained
        community.
\end{itemize}

\section{Conclusion}

We introduced LitJev, a trained single-pass decision model for manuscript
screening. It provides twelve independent calibrated decision heads and an
explicit coverage signal. On an in-domain held-out benchmark it is strong on
craft, integrity, and coverage tasks, while its rare value-class tasks remain
data-limited. The model is therefore positioned as a rapid screening and
calibration instrument, not as a replacement for expert review.

\section*{Acknowledgements}

We thank Dongbi Technology Data (Dongbi Data), the Chinese Academy of Sciences,
and Professor Dengsheng Wu of Shenzhen University for foundational
philosophical contributions to the principles of paper evaluation.

\section*{References}

\begin{thebibliography}{9}
\bibitem{jev} TypeSafe AI. Jev: typed, calibrated decisions.
\bibitem{openjev} OpenJev / mini-jev. Open reimplementation of a System-1 decision model.
\bibitem{rlcr} RLCR. Reinforcement Learning for Calibrated Decisions. arXiv:2507.16806.
\bibitem{reviewguard} ReviewGuard. Aligning review models with future citation outcomes. arXiv:2606.24892.
\bibitem{apres} APRES. Rubric discovery for predictable future citations. arXiv:2603.03142.
\bibitem{citation} Citation prediction is more learnable than review scores. arXiv:2503.05712.
\bibitem{kml} Yu Yang. Knowledge Manifold Learning, 2026.
\bibitem{qgt} QGT V2/V3. Preceding manifold-quality studies in the same project repository.
\end{thebibliography}

\end{document}
